\documentclass[11pt,a4paper]{article}
\usepackage[T1]{fontenc}
\usepackage{lmodern}
\usepackage{amsmath,amssymb}
\usepackage{booktabs}
\usepackage{array}
\usepackage[a4paper,margin=2.5cm]{geometry}
\usepackage[hidelinks]{hyperref}
\usepackage{amsthm}

\newif\ifshowworked \showworkedtrue
\newenvironment{worked}{\ifshowworked\par\smallskip\noindent\textit{Working.}\par\else\setbox0\vbox\bgroup\fi}{\ifshowworked\par\smallskip\else\egroup\fi}

\newcommand{\Nc}{\mathcal{N}}
\newcommand{\Gc}{\mathcal{G}}
\newcommand{\eps}{\varepsilon}
\newcommand{\bth}{\beta_{\mathrm{th}}}
\newcommand{\Pth}{\Pi_{\mathrm{th}}}
\newcommand{\ii}{\mathrm{i}}
\newcommand{\dd}{\mathrm{d}}
\newcommand{\bx}{\hat u_x}
\newcommand{\bz}{\hat u_z}
\newcommand{\bp}{\hat\varphi}
\newcommand{\bn}{\widehat{\Delta n}}
\newcommand{\bt}{\hat\theta}
\newcommand{\kB}{k_{\mathrm{th}}}

\title{\bfseries Lamb waves in a piezoelectric semiconductor plate with\\[2pt]
Moore--Gibson--Thompson heat conduction and the nonlocal\\[2pt]
strain gradient size effect: the complete hand calculation}
\author{Dr.\ Vipin\\ Gurugram University}
\date{30 October 2026}

\begin{document}
\maketitle

\begin{abstract}
\noindent
We reduce the five coupled field equations of a piezoelectric semiconductor plate to a plane-wave problem
in the thickness coordinate. The plate has thickness $2h$, free faces at $z=\pm h$, and carries a wave
$\exp[\ii(kx-\omega t)]$ with complex $k$. Heat conduction is of Moore--Gibson--Thompson type, and the
elastic stress carries the nonlocal strain gradient operator of Lim, Zhang and Reddy. The calculation is
written out in full: constitutive components for the hexagonal class $6mm$, the five governing equations in
components, the $5\times5$ amplitude matrix, the degree of its determinant, the parity of that determinant,
the amplitude ratios from the cofactors, the seven boundary conditions per face, the matrix whose
determinant is the dispersion relation, and the hand limits. Every step is inside a block that can be
switched off; every result stays outside.
\end{abstract}

\noindent\textit{Keywords.} Lamb wave; piezoelectric semiconductor; Moore--Gibson--Thompson heat conduction;
nonlocal strain gradient; plate; dispersion.

\section{Introduction}

A Lamb wave is an elastic wave trapped in a plate. The two free faces reflect the wave back and forth, so the
plate carries a finite number of wave shapes across its thickness instead of a continuous spectrum. Each
shape travels along the plate with its own speed, and that speed depends on the product of frequency and
thickness. This is why a plate is useful: the same plate behaves as a thin beam at low frequency and as a
half space at high frequency, and the changeover carries information about the material.

Add a piezoelectric semiconductor and two things enter. First, the electric field is not imposed from
outside; it is generated inside by the strain and it acts back on the strain through the piezoelectric
constants. Second, the crystal carries mobile carriers whose density is not fixed. A strain pulse produces a
piezoelectric field, that field pushes the carriers, the carriers pile up or drain away, and their
redistribution changes both the electric field and the temperature. The wave therefore sees a material
with a slow internal relaxation, and its attenuation carries the signature of that relaxation. This is the
physics behind acoustic charge transport devices, and it is why the coupled problem is worth solving
carefully.

Heat in such a plate does not obey the classical diffusion equation when the frequency is high. The
classical equation predicts an infinite speed for a temperature signal, which is harmless at low frequency
and wrong at high frequency. The Moore--Gibson--Thompson (MGT) equation repairs this by keeping a relaxation
term of the heat flux and a second time derivative; it is the version that follows from a well posed
thermoelastic theory and it reduces to the classical equation when the two time constants are set to zero
\cite{Quintanilla2019}. We keep it because the frequencies at which the carrier and thermal relaxations
compete with the elastic wave lie in the same window.

The second size effect is the nonlocal strain gradient theory (NSGT). Two lengths enter. The gradient length
$l_g$ (metre) measures how strongly the strain gradient stiffens the material; a small sample with a steep
strain gradient stores extra energy, like a rope that resists bending more when the bend is sharp. The
nonlocal length $l_{\rm nl}$ (metre) measures how far the stress at a point is averaged over its
neighbourhood; averaging softens the response, the way a crowd absorbs the push of one person. NSGT carries
both, and the ratio $l_g/l_{\rm nl}$ decides whether a thin plate is stiffer or softer than the classical
one \cite{Barretta2018}. The differential form we use is the one Lim, Zhang and Reddy gave; we did not read
their paper, and we say below exactly what we read instead.

Waves in piezoelectric semiconductor plates have a long history. Drift-diffusion coupling with the elastic
field in a plate was set out by Yang and Zhou \cite{YangZhou2005}; their equations are the starting point
for the carrier part of our model, though we did not read their numbered equations and say so. On the
thermoelectric side, with Seebeck and Peltier effects kept in a continuum theory, the reference is Qu, Pan,
Zhu, Jin and Roy \cite{Qu2023}; we read the abstract and the stated scope, not the equations. Gupta and
co-workers used the MGT framework on a piezoelectric semiconductor with a void \cite{Gupta2025}; title and
abstract tell us that the combination exists, and we did not read their heat equation. What is missing in that literature is a single worked chain from the three-dimensional
equations to the boundary matrix of a plate, with the thermal and carrier relaxations kept and the two size
lengths kept. That chain is what this file gives.

The novelty we claim is narrow, and we state it in the words agreed with the supervisor. Carrier and thermal
fields are kept as full dynamic fields, not as a quasi-static correction. Heat conduction keeps the MGT
operator with one time derivative cancelled, as fixed in the model document. Both NSGT lengths are kept,
with the gradient operator on the elastic stress and the nonlocal operator on the inertia. Boundary
conditions carry the surface gradient term and the double traction of the gradient theory. No claim is made
about a new material, a new numerical method, or a new frequency range.

\section{Setting and assumptions}

The plate occupies $-h\le z\le h$ and is unbounded in $x$ and $y$. Its faces $z=\pm h$ are free of
traction, free of surface charge, free of carrier flux and thermally insulating. The material is a
hexagonal piezoelectric semiconductor of class $6mm$ with the six-fold axis along $z$.

We take plane strain in the $x$--$z$ plane: $u_y=0$ and $\partial_y=0$ for every field. Nothing depends on
$y$. The wave travels along $x$ with the factor $\exp[\ii(kx-\omega t)]$, where $\omega$ (rad/s) is real and
$k=k_r+\ii k_i$ (rad/m) is complex. With this sign convention a positive $k_i$ means decay along $+x$; the
time factor is $\exp[-\ii\omega t]$.

Every field through the thickness is a sum of exponentials $\exp(\eta z)$, with $\eta$ (1/m) complex. We do
not assume the parity of $\eta$ in advance; Section~\ref{sec:D} shows when the roots come in pairs.

The theory is linear. Strains, electric fields, carrier perturbation $\Delta n$ (m$^{-3}$) and temperature
change $\theta$ (K) are all small, and every product of two small quantities is dropped. The carriers are of
one type, electrons, with equilibrium density $n_0$ (m$^{-3}$) uniform in the reference state. The
equilibrium electric field is zero. Recombination is linear with time constant $\tau_R$ (s), and it may be
switched off by letting $\tau_R\to\infty$.

Two operators carry the size effect. They are $\Nc(\eta)=1-l_{\rm nl}^2(\eta^2-k^2)$ and
$\Gc(\eta)=1-l_g^2(\eta^2-k^2)$; the combination $\eta^2-k^2$ is the Laplacian acting on a plane wave in
our sign convention, and the sign of each operator is fixed by the definition in Section~\ref{sec:src}.
The elastic stress is local in the nonlocal length and nonlocal in the gradient length only in the sense of
the differential law; the momentum balance carries the nonlocal operator on the inertia, which is the
equivalent differential form of the integral law.

\section{Notation}\label{sec:not}

\begin{center}\footnotesize
\begin{tabular}{p{2.1cm}p{5.6cm}p{1.9cm}p{4.6cm}}
\toprule
symbol & meaning in words & unit & where it enters\\
\midrule
$h$ & half the plate thickness; the faces sit at $\pm h$ & m & boundary rows\\
$k$ & complex wave number along the plate; its real part sets the wavelength & rad/m & every entry\\
$\omega$ & angular frequency, real and positive & rad/s & every entry, through $\ii\omega$\\
$\eta$ & exponential rate of a partial wave across the thickness & 1/m & roots of $\det\mathbf M$\\
$u_x,u_z$ & displacements along the plate and across it & m & rows 1 and 2\\
$\varphi$ & electric potential; the field is minus its gradient & V & row 3\\
$\Delta n$ & excess electron density over the equilibrium value & m$^{-3}$ & row 4\\
$\theta$ & change of temperature from the reference value & K & row 5\\
$T_0$ & reference temperature & K & heat row\\
$c_{11},c_{13},c_{33},c_{44}$ & elastic stiffnesses for the hexagonal class & Pa & rows 1, 2\\
$e_{31},e_{33},e_{15}$ & piezoelectric constants; stress per unit field & C/m$^2$ & rows 1--3\\
$\epsilon_{11},\epsilon_{33}$ & dielectric permittivities at constant strain & F/m & row 3\\
$\beta_1,\beta_3$ & thermal stress coefficients, $\beta_{ij}=c_{ijkl}\alpha_{kl}$ & Pa/K & rows 1, 2, 5\\
$p_3$ & pyroelectric constant; charge released per kelvin & C/(m$^2$K) & rows 3, 5\\
$\alpha_{ij}$ & thermal expansion coefficients & 1/K & through $\beta$\\
$q$ & elementary charge, a positive number & C & row 3\\
$n_0$ & equilibrium electron density & m$^{-3}$ & row 4\\
$b_n$ & electron mobility & m$^2$/(V\,s) & row 4\\
$D_n$ & electron diffusivity & m$^2$/s & row 4\\
$S_n$ & Seebeck coefficient of the electrons, negative for $n$-type & V/K & rows 4, 5\\
$v_d$ & drift velocity along the plate & m/s & row 4\\
$\bth$ & thermodiffusion coefficient, $\bth=-b_nn_0S_n$, positive & A/(m\,K) & row 4\\
$\Pth$ & Peltier coefficient, $\Pth=T_0S_n$, negative for $n$-type & V & row 5\\
$\tau_R$ & recombination time; infinite when recombination is off & s & row 4\\
$E_g$ & band gap & J & row 5\\
$\rho$ & mass density & kg/m$^3$ & rows 1, 2\\
$c_p$ & specific heat at constant strain & J/(kg\,K) & row 5\\
$\kB$ & thermal conductivity & W/(m\,K) & row 5\\
$\tau_q$ & MGT relaxation time of the heat flux & s & row 5\\
$\kappa$ & thermal diffusivity, $\kB/(\rho c_p)$ & m$^2$/s & row 5\\
$l_g$ & strain gradient length of NSGT & m & $\Gc$\\
$l_{\rm nl}$ & nonlocal length of NSGT & m & $\Nc$\\
$\mathbf M(\eta)$ & $5\times5$ amplitude matrix of one partial wave & --- & Section~\ref{sec:C}\\
$\mathbf B$ & boundary matrix, one row per condition, one column per root & --- & Section~\ref{sec:G}\\
\bottomrule
\end{tabular}
\end{center}

One everyday picture helps before the algebra. Think of the plate as a corridor with two hard walls. A
push travels along the corridor and bounces between the walls; only those pushes whose bounce pattern fits
the corridor exactly survive, and each surviving pattern is a mode. The algebra below is only the
bookkeeping of which bounce patterns fit.

\section{Source equations}\label{sec:src}

We list the equations we start from, one at a time, with what we read and how we use it. The tag at the end
of each item is one of three: taken from the source as it stands, taken from the source with the change we
name, or ours because it follows from a named law.

\subsection{Strain and field}

The strain is the symmetric part of the displacement gradient, $\eps_{ij}=\tfrac12(u_{i,j}+u_{j,i})$, and
the electric field is minus the gradient of the potential, $E_i=-\varphi_{,i}$. These are textbook laws of
linear elasticity and electrostatics; no equation number is needed and none is claimed. Under our plane
wave they give $\eps_{xx}=\ii k\,\bx$, $\eps_{zz}=\eta\,\bz$,
$\eps_{xz}=\tfrac12(\eta\bx+\ii k\bz)$, $E_x=-\ii k\bp$, $E_z=-\eta\bp$. 
Tag: textbook law.

\subsection{Constitutive relations of the piezoelectric semiconductor}

For class $6mm$ with the axis along $z$, the stress and the electric displacement are
\begin{equation}\label{eq:const}
\begin{aligned}
\sigma_{ij}&=c_{ijkl}\bigl(1-l_g^2\nabla^2\bigr)\eps_{kl}-e_{kij}E_k-\beta_{ij}\theta,\\
D_i&=e_{ikl}\eps_{kl}+\epsilon_{ik}E_k+p_i\theta .
\end{aligned}
\end{equation}
The first line is the piezoelectric constitutive law of Zhu, Negahban, Xu, Xia and Li
\cite[Eqs.~(1)--(3)]{Zhu2023} with two additions: the thermal stress $-\beta_{ij}\theta$, which
is Duhamel's law and carries no equation number, and the gradient operator on the elastic strain, which is
the NSGT placement explained below. The second line is the same source \cite[Eqs.~(1)--(3)]{Zhu2023} plus the pyroelectric term
$p_i\theta$, a textbook term.
Tag: taken from the source for the piezoelastic part; the gradient operator on the elastic strain and the
two thermal terms are ours and follow from the named laws.

We must be precise about what we read in \cite{Zhu2023}. Their Eqs.~(1) to (7) were read as a block, and
that block contains the strain--displacement relation, the piezoelectric constitutive relations, Gauss's law
with a mobile charge density, the drift--diffusion current and the continuity equation for the carriers. No
note was kept of which number belongs to which line. Before submission, please open that block and write
the exact number beside each of our equations; until you do, the citation you see here is to the block only.

\subsection{Balance of momentum with the nonlocal operator}

The momentum balance is
\begin{equation}\label{eq:mom}
\bigl(1-l_{\rm nl}^2\nabla^2\bigr)\rho\,\ddot u_i=\sigma_{ij,j}.
\end{equation}
The differential form is the one quoted in the introduction of Barretta
and Marotti de Sciarra \cite[Eqs.~(18)--(19)]{Barretta2018}, who attribute it to Lim, Zhang and Reddy
\cite{Lim2015}. We did not open \cite{Lim2015}; the equations (17) to (19) of that paper are known to us
only through \cite[Eqs.~(18)--(19)]{Barretta2018}. Section~2 of Lim et al.\ is the place to confirm the
differential law and the constitutive boundary conditions, and you should open it before you defend the
placement.
Tag: taken from \cite{Barretta2018} as a differential law; the placement of the nonlocal operator on the
inertia alone is ours, and it is the equivalent differential form of the integral law.

\subsection{Gauss's law with mobile charge}

Gauss's law reads
\begin{equation}\label{eq:gauss}
D_{i,i}=-q\,\Delta n .
\end{equation}
It is inside the block \cite[Eqs.~(1)--(7)]{Zhu2023}. The right side
is the charge of the excess electrons; the sign is fixed by the convention that electrons carry charge $-q$
with $q>0$, so an excess of electrons is a negative charge density.
Tag: taken from the source.

\subsection{Carrier continuity}

The carrier balance is
\begin{equation}\label{eq:carr}
\frac{\partial\Delta n}{\partial t}
=D_n\nabla^2\Delta n+b_nn_0\,\nabla\!\cdot\!E
+v_d\frac{\partial\Delta n}{\partial x}+\bth\nabla^2\theta-\frac{\Delta n}{\tau_R}.
\end{equation}
The first two terms on the right are drift and diffusion as written in
the block \cite[Eqs.~(1)--(7)]{Zhu2023}; the drift term with $v_d$ is the same drift written in the frame
moving with the carriers; the last two terms need their own justification.
The thermodiffusion term $\bth\nabla^2\theta$ follows from irreversible thermodynamics: a temperature
gradient drives a particle current, with the Onsager cross coefficient fixed by the Seebeck coefficient as
$\bth=-b_nn_0S_n$. With $S_n<0$ for electrons this coefficient is positive, so carriers are driven towards
the cold side, which is the behaviour seen in thermoelectric materials.
The recombination term $-\Delta n/\tau_R$ is the linearised decay of the excess density.
Tag: the drift and diffusion terms are taken from the source; the thermodiffusion term and the
recombination term are ours and follow from the Onsager reciprocity and the linearised decay law.

\subsection{MGT heat equation}

The heat equation is
\begin{equation}\label{eq:heat}
\Bigl(1+\tau_q\frac{\partial}{\partial t}\Bigr)
\Bigl[\rho c_p\dot\theta+T_0\beta_{ij}\dot\eps_{ij}+T_0p_i\dot E_i
+q\Pth\frac{\partial\Delta n}{\partial t}+\frac{E_g\Delta n}{\tau_R}\Bigr]
=\kB\nabla^2\theta .
\end{equation}
The operator on the left is the Moore--Gibson--Thompson operator with one
time derivative cancelled, which is the form read in \cite[Eqs.~(2.5) and (3.1)--(3.5)]{Quintanilla2019}; in
that form the heat flux has a relaxation time $\tau_q$ and the temperature has the parabolic principal part.
We use it with the thermal stress coefficients $\beta_1,\beta_3$ and the pyroelectric constant $p_3$. The
last two terms inside the bracket are the Peltier heat of the carrier current and the recombination heat,
and they are the partners of the thermodiffusion and recombination terms in \eqref{eq:carr}.
Tag: the operator is taken from \cite{Quintanilla2019}; the coupling bracket is the classical
thermoelastic one, with the pyroelectric and Peltier terms added as ours from the stated laws.

The two cross coefficients are not independent. With $\bth=-b_nn_0S_n$ and $\Pth=T_0S_n$,
\begin{equation}\label{eq:onsager}
\frac{\Pth}{\bth}=-\frac{T_0}{b_nn_0},
\end{equation}
This is the Onsager relation $\Pi=TS$ of textbook thermoelectricity. For
$n$-type electrons $S_n<0$, so $\bth>0$ and $\Pth<0$. In \eqref{eq:heat} a positive
$\partial_t\Delta n$ with $\Pth<0$ subtracts heat, which is what physics asks: carriers that pile up carry
their heat with them. 
Tag: textbook thermoelectricity, cited without an equation number.

\subsection{Why the heat equation has this shape}

The form of \eqref{eq:heat} is not a guess. It follows from two statements: an energy balance and a flux
law. The energy balance for the small temperature change is
\begin{equation}\label{eq:energy}
\rho c_p\dot\theta+T_0\beta_{ij}\dot\eps_{ij}+T_0p_i\dot E_i
+q\Pth\frac{\partial\Delta n}{\partial t}+\frac{E_g\Delta n}{\tau_R}=-q_{i,i},
\end{equation}
where $q_i$ (W/m$^2$) is the heat flux. Each term on the left is a heat source per unit volume: the first is
the heat stored, the second is the heat converted to elastic stress by thermal expansion, the third is the
heat absorbed by the pyroelectric polarisation as the field changes, the fourth is the Peltier heat carried
by the carriers, and the fifth is the heat released when carriers recombine. The MGT flux law, with one time
derivative cancelled, is
\begin{equation}\label{eq:flux}
q_i+\tau_q\dot q_i=-\kB\theta_{,i}.
\end{equation}

\begin{worked}
Take the divergence of \eqref{eq:flux}: $q_{i,i}+\tau_q\dot q_{i,i}=-\kB\nabla^2\theta$. From
\eqref{eq:energy}, $q_{i,i}=-B$, where $B$ is the left side of \eqref{eq:energy} without its last term. Then
$\dot q_{i,i}=-\dot B$, and the diverged flux law reads $-B-\tau_q\dot B=-\kB\nabla^2\theta$, that is
$(1+\tau_q\partial_t)B=\kB\nabla^2\theta$. Writing $B$ out gives \eqref{eq:heat} with $f=1-\ii\omega\tau_q$
under the plane wave. The sign is worth one check: with $\tau_q=0$ the equation returns the classical
Fourier form, and with $\theta$ uniform and no sources it returns $\rho c_p\dot\theta=0$, as it should.
\end{worked}

\noindent This derivation also tells you what the two time constants mean. In the flux law \eqref{eq:flux}
the flux cannot change instantaneously, and the relaxation time $\tau_q$ is the delay. A second time constant,
$\tau_T$ of the full MGT theory, appears when the energy balance itself is written with a lag; we drop it,
which is the standard weak-damping form, and the effect of keeping it is one extra term in $M_{55}$.

The two cross coefficients of \eqref{eq:carr} and \eqref{eq:heat} must be equal by Onsager symmetry. Write
the entropy production of the coupled carrier and heat fluxes; it is a bilinear form in the temperature
gradient and the electrochemical potential gradient, so the coefficient matrix must be symmetric. One entry
of that matrix is the thermodiffusion coefficient of the carrier equation, the other is the Peltier
coefficient of the heat equation, and their ratio is fixed by the equilibrium carrier density and the
mobility. Equation \eqref{eq:onsager} is that statement written for electrons, and it is the only line
in the model you cannot check by algebra alone: it needs the entropy production written out.

\subsection{The MGT time constants}

The phase lag $\tau_T$ of the MGT equation is set by the thickness through $\tau_T=(2h)^2/(\pi^2\kappa)$,
the time a temperature disturbance needs to cross the plate once. The equation we use keeps $\tau_q$ and
drops the term with $\tau_T^2$, which is the standard weak-damping form. If you want the full three-time
form, it is a one-line change in one entry of the matrix and nothing else changes; we state it in
Section~\ref{sec:H}.

\section{Calculation A: constitutive components for class $6mm$}\label{sec:A}

The point of this section is to say which material constants survive in our geometry, and why the rest
drop out. We do it in words and then in components.

\subsection{Which components survive}

\begin{worked}
The crystal has a six-fold axis along $z$. A rotation of $60^\circ$ about $z$ must leave the constants
unchanged. Write the constants in matrix notation with $xx\to1$, $zz\to3$, $xz\to5$. Under the rotation the
in-plane shear $4$ goes to a combination of $4$ and $5$, and $5$ goes to a combination of $5$ and $4$; the
only way both stay invariant is $c_{45}=0$ together with the two in-plane shears equal, which is
$c_{66}=(c_{11}-c_{12})/2$. The shear on the $x$--$z$ plane is $c_{44}$ and on the $y$--$z$ plane is $c_{55}$;
the rotation exchanges them, so $c_{44}=c_{55}$. A rotation of $180^\circ$ about $x$ or $y$ takes $z$ to
$-z$ and is not a symmetry of a poled crystal, so terms with an odd number of $z$ indices are allowed:
$e_{31}$, $e_{33}$, $e_{15}$ and $p_3$ survive.
\end{worked}

In our plane strain with no $y$ dependence the fields that survive are $\eps_{xx},\eps_{zz},\eps_{xz}$,
$E_x,E_z$ and $\theta$. Their list fixes the constants, and the list splits in two. For the three stresses
$\sigma_{xx},\sigma_{zz},\sigma_{xz}$ the constants are $c_{11},c_{13},c_{33},c_{44}$ together with the two
thermal stress coefficients $\beta_1,\beta_3$. For the two induction components $D_x,D_z$ they are
$e_{31},e_{33},e_{15}$, $\epsilon_{11},\epsilon_{33}$ (components of the $3\times3$ dielectric tensor, so
$\epsilon_{11}=\epsilon_{22}$ and $\epsilon_{33}$ is the axial one) and the pyroelectric constant $p_3$.
The constant $c_{12}$ never appears because $\eps_{yy}=0$; it would appear only through $c_{66}$ if we kept an in-plane
shear wave, and even then it enters through the combination $(c_{11}-c_{12})/2$.

\subsection{The stresses and the induction in components}

Putting the plane wave into \eqref{eq:const}, and writing $\Gc=1-l_g^2(\eta^2-k^2)$,
\begin{align}
\hat\sigma_{xx}&=\Gc\bigl(\ii k\,c_{11}\bx+\eta\,c_{13}\bz\bigr)+\eta\,e_{31}\bp-\beta_1\bt,
\label{eq:sxx}\\
\hat\sigma_{zz}&=\Gc\bigl(\ii k\,c_{13}\bx+\eta\,c_{33}\bz\bigr)+\eta\,e_{33}\bp-\beta_3\bt,
\label{eq:szz}\\
\hat\sigma_{xz}&=\Gc\,c_{44}\bigl(\eta\bx+\ii k\bz\bigr)+\ii k\,e_{15}\bp,
\label{eq:sxz}\\
\hat D_x&=e_{15}\bigl(\eta\bx+\ii k\bz\bigr)-\ii k\,\epsilon_{11}\bp,
\label{eq:Dx}\\
\hat D_z&=\ii k\,e_{31}\bx+\eta\,e_{33}\bz-\eta\,\epsilon_{33}\bp+p_3\bt .
\label{eq:Dz}
\end{align}

\begin{worked}
Each line is the same routine, so one line is enough to show it. In \eqref{eq:sxx} the term $c_{11}\eps_{xx}$
becomes $c_{11}(\ii k\bx)$ because the only derivative is $\partial_x\to\ii k$; the term $c_{13}\eps_{zz}$
becomes $c_{13}(\eta\bz)$; the piezoelectric term $-e_{kij}E_k$ contributes $-e_{31}E_z=-e_{31}(-\eta\bp)$,
which is $+\eta e_{31}\bp$; the thermal term is $-\beta_1\bt$. The gradient operator multiplies only the two
elastic terms, never the piezoelectric or thermal ones, because in \eqref{eq:const} it acts on $\eps_{kl}$
alone. The other four lines follow the same routine with the constants of the class: for $D_z$ the relevant
piezoelectric constant is $e_{33}$ acting on $\eps_{zz}$ and $e_{31}$ on $\eps_{xx}$, and the dielectric
term is $-\epsilon_{33}E_z=+\eta\epsilon_{33}\bp$; the pyroelectric term is $+p_3\bt$ with the sign of
\eqref{eq:const}.
\end{worked}

Two remarks on signs, because signs are where a first draft usually fails. A compressive $\theta>0$ with
$\beta>0$ reduces the tensile stress, which is why the thermal term is negative in \eqref{eq:sxx}; heating a
constrained plate puts it in compression. The pyroelectric term in \eqref{eq:Dz} is positive because a
positive $\theta$ releases bound charge of the same sign as the poling.

\section{Calculation B: the five governing equations in components}\label{sec:B}

The five unknown amplitudes are collected in a column,
\begin{equation}\label{eq:vec}
\mathbf a=\bigl(\bx,\ \bz,\ \bp,\ \bn,\ \bt\bigr)^{\!\top}.
\end{equation}
The five equations are the two components of \eqref{eq:mom}, Gauss \eqref{eq:gauss}, the carrier balance
\eqref{eq:carr} and the heat balance \eqref{eq:heat}.

\subsection{Momentum, $x$ component}

\begin{worked}
Take the $x$ component of \eqref{eq:mom}. The divergence is $\sigma_{xj,j}=\partial_x\sigma_{xx}+\partial_z\sigma_{xz}$,
so under the plane wave it becomes $\ii k\hat\sigma_{xx}+\eta\hat\sigma_{xz}$. The left side is
$\Nc\rho\,\ddot u_x$, and with $\ddot u_x=(-\ii\omega)^2\bx=-\omega^2\bx$ it is $-\Nc\rho\omega^2\bx$.
Bringing everything to one side,
\[
-\bigl(\ii k\hat\sigma_{xx}+\eta\hat\sigma_{xz}\bigr)-\Nc\rho\omega^2\bx=0 .
\]
\end{worked}

\noindent Result, first of the five:
\begin{equation}\label{eq:g1}
-\bigl(\ii k\hat\sigma_{xx}+\eta\hat\sigma_{xz}\bigr)-\Nc\rho\omega^2\bx=0 .
\end{equation}

\subsection{Momentum, $z$ component}

\begin{worked}
The $z$ component is $\sigma_{zj,j}=\partial_x\sigma_{xz}+\partial_z\sigma_{zz}\to\ii k\hat\sigma_{xz}+\eta\hat\sigma_{zz}$,
and the inertia term is again $-\Nc\rho\omega^2\bz$. Hence
\[
-\bigl(\ii k\hat\sigma_{xz}+\eta\hat\sigma_{zz}\bigr)-\Nc\rho\omega^2\bz=0 .
\]
\end{worked}

\noindent Result, second of the five:
\begin{equation}\label{eq:g2}
-\bigl(\ii k\hat\sigma_{xz}+\eta\hat\sigma_{zz}\bigr)-\Nc\rho\omega^2\bz=0 .
\end{equation}

\subsection{Gauss's law}

\begin{worked}
The divergence of the induction is $\partial_xD_x+\partial_zD_z$ and the charge is $-q\bn$. Equation
\eqref{eq:gauss} becomes
\[
\ii k\hat D_x+\eta\hat D_z=-q\bn ,
\]
and moving the right side across gives $\ii k\hat D_x+\eta\hat D_z+q\bn=0$.
\end{worked}

\noindent Result, third of the five:
\begin{equation}\label{eq:g3}
\ii k\hat D_x+\eta\hat D_z+q\bn=0 .
\end{equation}

\subsection{Carrier balance}

\begin{worked}
In \eqref{eq:carr} every time derivative is $-\ii\omega$, every $\partial_x$ is $\ii k$ and every Laplacian
is $\eta^2-k^2$. The left side is $-\ii\omega\bn$. The first right-hand term is $D_n(\eta^2-k^2)\bn$. The
field divergence in the second term is $\nabla\cdot E=-\nabla^2\varphi=-(\eta^2-k^2)\bp$, so that term is
$-b_nn_0(\eta^2-k^2)\bp$. The drift term is $v_d(\ii k\bn)$. The thermodiffusion term is
$\bth(\eta^2-k^2)\bt$, and the recombination term is $-\bn/\tau_R$. Bring all five to the left:
\[
-\ii\omega\bn-D_n(\eta^2-k^2)\bn+b_nn_0(\eta^2-k^2)\bp-\ii k v_d\bn-\bth(\eta^2-k^2)\bt+\frac{\bn}{\tau_R}=0 .
\]
\end{worked}

\noindent Result, fourth of the five:
\begin{equation}\label{eq:g4}
\Bigl[-\ii\omega-D_nL-\ii kv_d+\frac{1}{\tau_R}\Bigr]\bn+b_nn_0L\,\bp-\bth L\,\bt=0,
\qquad L=\eta^2-k^2 .
\end{equation}

\subsection{Heat balance}

\begin{worked}
Write $\mathrm{D}_t=-\ii\omega$ for the time derivative and $f=1+\tau_q\mathrm{D}_t=1-\ii\omega\tau_q$ for
the MGT operator. In the bracket of \eqref{eq:heat} the first term is $-\ii\omega\rho c_p\bt$. The thermal
stress term is $T_0\beta_{ij}\dot\eps_{ij}$, and in our plane strain the two terms that survive are
$T_0\beta_1\dot\eps_{xx}+T_0\beta_3\dot\eps_{zz}$, which give
$T_0\beta_1(-\ii\omega)(\ii k\bx)+T_0\beta_3(-\ii\omega)(\eta\bz)=T_0\omega k\beta_1\bx-\ii\omega T_0\beta_3\eta\bz$.
The pyroelectric term is $T_0p_i\dot E_i=T_0p_3\dot E_z=T_0p_3(-\ii\omega)(-\eta\bp)=\ii\omega T_0p_3\eta\bp$.
The Peltier term is $q\Pth(-\ii\omega)\bn$, and the recombination term is $E_g\bn/\tau_R$. Multiplying the
bracket by $f$ and subtracting the conduction term $\kB(\eta^2-k^2)\bt$ gives
\[
f\Bigl[-\ii\omega\rho c_p\bt+T_0\omega k\beta_1\bx-\ii\omega T_0\beta_3\eta\bz
+\ii\omega T_0p_3\eta\bp-\ii\omega q\Pth\bn+\frac{E_g}{\tau_R}\bn\Bigr]-\kB L\bt=0 .
\]
\end{worked}

\noindent Result, fifth of the five:
\begin{equation}\label{eq:g5}
f\Bigl[-\ii\omega\rho c_p\bt+T_0\omega k\beta_1\bx-\ii\omega T_0\beta_3\eta\bz
+\ii\omega T_0p_3\eta\bp+\Bigl(-\ii\omega q\Pth+\frac{E_g}{\tau_R}\Bigr)\bn\Bigr]-\kB L\bt=0 .
\end{equation}

\section{Calculation C: the amplitude matrix $\mathbf M(\eta)$}\label{sec:C}

The five equations \eqref{eq:g1}--\eqref{eq:g5} are linear in the five amplitudes and homogeneous. Read off
the coefficient of each amplitude in each equation, in the order of \eqref{eq:vec}. Entries not listed are
zero.

\subsection{Row 1, momentum along the plate}

\begin{worked}
Put \eqref{eq:sxx} and \eqref{eq:sxz} into \eqref{eq:g1} and collect.
For $\bx$: from $\ii k\hat\sigma_{xx}$ we get $\ii k\,\Gc c_{11}(\ii k)=-k^2\Gc c_{11}$; from
$\eta\hat\sigma_{xz}$ we get $\eta\,\Gc c_{44}\eta=\eta^2\Gc c_{44}$; the inertia contributes
$-\Nc\rho\omega^2$. With the overall minus in \eqref{eq:g1} the two elastic pieces change sign and the
inertia does not: $M_{11}=-\Nc\rho\omega^2+\Gc(c_{11}k^2-c_{44}\eta^2)$.
For $\bz$: $\ii k\cdot\Gc c_{13}\eta$ and $\eta\cdot\Gc c_{44}(\ii k)$ add to $\ii k\eta\Gc(c_{13}+c_{44})$,
and the overall minus gives $M_{12}=-\ii k\eta\Gc(c_{13}+c_{44})$.
For $\bp$: $\ii k\cdot\eta e_{31}$ and $\eta\cdot\ii k e_{15}$ add to $\ii k\eta(e_{31}+e_{15})$, with a minus:
$M_{13}=-\ii k\eta(e_{31}+e_{15})$.
For $\bn$: nothing, since neither stress contains $\bn$: $M_{14}=0$.
For $\bt$: $\ii k\cdot(-\beta_1)$, with the overall minus, gives $M_{15}=\ii k\beta_1$.
\end{worked}

\noindent The first row is
\begin{equation}\label{eq:row1}
M_{11}=-\Nc\rho\omega^2+\Gc\bigl(c_{11}k^2-c_{44}\eta^2\bigr),\quad
M_{12}=-\ii k\eta\,\Gc(c_{13}+c_{44}),\quad
M_{13}=-\ii k\eta(e_{31}+e_{15}),
\end{equation}
\begin{equation}\label{eq:row1b}
M_{14}=0,\qquad M_{15}=\ii k\beta_1 .
\end{equation}

\subsection{Row 2, momentum across the plate}

\begin{worked}
Same routine with \eqref{eq:g2}. For $\bx$: $\ii k\cdot\Gc c_{44}\eta$ plus $\eta\cdot\Gc c_{13}(\ii k)$ gives
$\ii k\eta\Gc(c_{13}+c_{44})$, with a minus.
For $\bz$: $\ii k\cdot\Gc c_{44}(\ii k)=-k^2\Gc c_{44}$ and $\eta\cdot\Gc c_{33}\eta=\eta^2\Gc c_{33}$; with the
overall minus and the inertia term, $M_{22}=-\Nc\rho\omega^2+\Gc(c_{44}k^2-c_{33}\eta^2)$.
For $\bp$: $\ii k\cdot\ii k e_{15}=-k^2e_{15}$ and $\eta\cdot\eta e_{33}=\eta^2e_{33}$, with a minus:
$M_{23}=e_{15}k^2-e_{33}\eta^2$.
For $\bt$: $\eta\cdot(-\beta_3)$ with the minus gives $M_{25}=\eta\beta_3$. The $\bn$ entry is zero.
\end{worked}

\noindent The second row is
\begin{equation}\label{eq:row2}
M_{21}=-\ii k\eta\,\Gc(c_{13}+c_{44}),\quad
M_{22}=-\Nc\rho\omega^2+\Gc\bigl(c_{44}k^2-c_{33}\eta^2\bigr),\quad
M_{23}=e_{15}k^2-e_{33}\eta^2,
\end{equation}
\begin{equation}\label{eq:row2b}
M_{24}=0,\qquad M_{25}=\eta\beta_3 .
\end{equation}

\subsection{Row 3, Gauss's law}

\begin{worked}
Use \eqref{eq:Dx} and \eqref{eq:Dz} in \eqref{eq:g3}. For $\bx$: $\ii k\cdot e_{15}\eta+\eta\cdot\ii k e_{31}$,
so $M_{31}=\ii k\eta(e_{15}+e_{31})$.
For $\bz$: $\ii k\cdot e_{15}(\ii k)=-k^2e_{15}$ and $\eta\cdot\eta e_{33}=\eta^2e_{33}$, so
$M_{32}=e_{33}\eta^2-e_{15}k^2$.
For $\bp$: $\ii k\cdot(-\ii k\epsilon_{11})=k^2\epsilon_{11}$ and $\eta\cdot(-\eta\epsilon_{33})=-\eta^2\epsilon_{33}$,
so $M_{33}=\epsilon_{11}k^2-\epsilon_{33}\eta^2$.
For $\bn$: the term $+q\bn$ in \eqref{eq:g3} gives $M_{34}=q$.
For $\bt$: $\eta\cdot p_3$ gives $M_{35}=\eta p_3$.
\end{worked}

\noindent The third row is
\begin{equation}\label{eq:row3}
M_{31}=\ii k\eta(e_{15}+e_{31}),\quad
M_{32}=e_{33}\eta^2-e_{15}k^2,\quad
M_{33}=\epsilon_{11}k^2-\epsilon_{33}\eta^2,\quad
M_{34}=q,\quad
M_{35}=\eta p_3 .
\end{equation}

\subsection{Row 4, carrier balance}

\begin{worked}
Equation \eqref{eq:g4} is already in matrix form. The three nonzero entries are the coefficient of $\bn$,
that of $\bp$ and that of $\bt$: $M_{44}=-\ii\omega-D_nL-\ii kv_d+1/\tau_R$, $M_{43}=b_nn_0L$ and
$M_{45}=-\bth L$. The entries $M_{41}$ and $M_{42}$ are zero because the carrier equation feels the
displacement only through the field, and the field only through $\bp$.
\end{worked}

\noindent The fourth row is
\begin{equation}\label{eq:row4}
M_{41}=0,\quad M_{42}=0,\quad M_{43}=b_nn_0L,\quad
M_{44}=-\ii\omega-D_nL-\ii kv_d+\frac{1}{\tau_R},\quad M_{45}=-\bth L .
\end{equation}

\subsection{Row 5, heat balance}

\begin{worked}
Read the coefficients from \eqref{eq:g5}. For $\bx$: $fT_0\omega k\beta_1$. For $\bz$:
$-f\ii\omega T_0\beta_3\eta$. For $\bp$: $f\ii\omega T_0p_3\eta$. For $\bn$:
$f(-\ii\omega q\Pth+E_g/\tau_R)$. For $\bt$: $f(-\ii\omega\rho c_p)-\kB L$.
\end{worked}

\noindent The fifth row is
\begin{equation}\label{eq:row5}
M_{51}=fT_0\omega k\beta_1,\quad
M_{52}=-\ii\omega fT_0\beta_3\eta,\quad
M_{53}=\ii\omega fT_0p_3\eta,\quad
M_{54}=f\Bigl(-\ii\omega q\Pth+\frac{E_g}{\tau_R}\Bigr),\quad
M_{55}=-\ii\omega f\rho c_p-\kB L .
\end{equation}

\subsection{The matrix and the two operators}

Collecting \eqref{eq:row1}--\eqref{eq:row5}, the five equations read
\begin{equation}\label{eq:M}
\mathbf M(\eta)\,\mathbf a=\mathbf 0,\qquad
\mathbf a=(\bx,\bz,\bp,\bn,\bt)^{\top},
\end{equation}
with the rows in the order momentum $x$, momentum $z$, Gauss, carriers, heat, and with
\begin{equation}\label{eq:NG}
\Nc(\eta)=1-l_{\rm nl}^2(\eta^2-k^2),\qquad
\Gc(\eta)=1-l_g^2(\eta^2-k^2),\qquad
L=\eta^2-k^2,\qquad
f=1-\ii\omega\tau_q .
\end{equation}
The matrix is written as five rows because a $5\times5$ array with long entries runs off the page. When you
copy it into your own notes, keep the row order; every later formula in this file assumes this order.

Notice that $\Nc$ appears in two places only, $M_{11}$ and $M_{22}$, always on the inertia; that $\Gc$
appears in the six elastic entries of rows 1 and 2 and nowhere else; and that the electric, carrier and
thermal entries carry no size operator at all. That is the content of the placement decision: the size
effect enters through the two operators and through the two boundary conditions of Section~\ref{sec:F}, and
nowhere else.

\section{Calculation D: the characteristic equation}\label{sec:D}

Nonzero amplitudes exist only if the determinant vanishes,
\begin{equation}\label{eq:char}
\det\mathbf M(\eta)=0 .
\end{equation}
This is a polynomial equation in $\eta$, and the number of its roots is the number of partial waves we must
superpose. So we need its degree, and we need to know whether the roots come in pairs.

\subsection{Degree counting}

\begin{worked}
Go through the matrix and write the highest power of $\eta$ in each row, allowing for the fact that the
permutation in the determinant may take any entry of a row.
Row 1: $M_{11}$ has $\Gc\eta^2\sim\eta^4$ through the product $l_g^2\eta^2\cdot c_{44}\eta^2$; $M_{12}$ has
$\Gc\eta\sim\eta^3$; $M_{13}\sim\eta$; $M_{15}\sim\eta^0$.
Row 2: $M_{21}\sim\eta^3$, $M_{22}\sim\eta^4$, $M_{23}\sim\eta^2$, $M_{25}\sim\eta^1$.
Row 3: $M_{31}\sim\eta^1$, $M_{32}\sim\eta^2$, $M_{33}\sim\eta^2$, $M_{34}\sim\eta^0$, $M_{35}\sim\eta^1$.
Row 4: $M_{43}\sim\eta^2$, $M_{44}\sim\eta^2$, $M_{45}\sim\eta^2$, others zero.
Row 5: $M_{51}\sim\eta^0$, $M_{52}\sim\eta^1$, $M_{53}\sim\eta^1$, $M_{54}\sim\eta^0$, $M_{55}\sim\eta^2$.
A determinant is a sum of products, one entry from each row and each column, so its degree is at most the
largest total $\sum_i d_i$ over all permutations. The permutation that reaches the maximum takes
$M_{11}$ (degree 4), $M_{22}$ (4), $M_{33}$ (2), $M_{44}$ (2) and $M_{55}$ (2), which sums to 14. No other
permutation gives more, because rows 1 and 2 can give at most 4 each and rows 3, 4, 5 at most 2 each. So the
degree is 14 when $l_g\ne0$.
\end{worked}

\noindent When $l_g=0$ the operator $\Gc$ is the number one, and the elastic entries of rows 1 and 2 drop
from degree 4 and 3 to degree 2 and 1. The same counting gives
\begin{equation}\label{eq:deg}
\deg_\eta\det\mathbf M=
\begin{cases}
14, & l_g\ne0,\\
10, & l_g=0 .
\end{cases}
\end{equation}
The nonlocal length does not change the degree: $\Nc$ multiplies only the constant inertia terms in
$M_{11}$ and $M_{22}$, so it adds powers of $\eta^2$ to entries that are already of higher degree in the
same positions. That is why $l_{\rm nl}$ changes the roots without changing their number.

\subsection{Parity of the determinant}

Under $\eta\to-\eta$ an entry changes sign if and only if it contains an odd power of $\eta$. We ask whether
a diagonal sign change of the five amplitudes can absorb all those sign changes, because if it can, then
the determinant is even in $\eta$ and the roots come in $\pm$ pairs.

\begin{worked}
Write $p(i,j)$ for the parity of the entry in row $i$, column $j$, with $E$ for even and $O$ for odd, and
read the parities from \eqref{eq:row1}--\eqref{eq:row5}:
row 1 is $E,O,O,E$ for the columns $\bx,\bz,\bp,\bt$ with $M_{14}=0$; row 2 is $O,E,E,O$; row 3 is $O,E,E,O$
for $\bx,\bz,\bp,\bt$ with $M_{34}=q$ even; row 4 is $E,E,E$ in the columns $\bp,\bn,\bt$; row 5 is
$E,O,O,E,E$ across all five columns.
Suppose a sign vector $s=(s_1,\dots,s_5)$ exists with $p(i,j)=s_is_j$. Take $s_1=-1$, which is free because
$M_{11}$ is even. Then $M_{12}$ odd gives $s_2=+1$; $M_{13}$ odd gives $s_3=+1$; $M_{15}$ even gives
$s_5=-1$; $M_{34}$ even gives $s_4=s_3=+1$. Now check the remaining entries: $M_{23}$ even needs
$s_2s_3=+1$, and it is $+1$, so fine; $M_{25}$ odd needs $s_2s_5=-1$, and it is $-1$, fine; $M_{35}$ odd
needs $s_3s_5=-1$, fine. The one entry that fails is $M_{45}=-\bth L$, which is even but needs
$s_4s_5=-1$. The same failure comes from $M_{54}$ when $\Pth\ne0$, since it is even and needs $s_4s_5=-1$ as
well.
\end{worked}

\noindent Hence
\begin{equation}\label{eq:parity}
\det\mathbf M(-\eta)=\det\mathbf M(\eta)
\quad\Longleftrightarrow\quad
\bth=0\ \text{and}\ \Pth=0 ,
\end{equation}
and in that case the sign vector is $(-1,+1,+1,+1,-1)$ in the order of \eqref{eq:vec}, that is
$\bx\to-\bx$ and $\bt\to-\bt$ with the other three unchanged. With either thermoelectric coefficient present
no such vector exists, the determinant is neither even nor odd, and the roots do not pair.

The physical reading is short. A temperature gradient that drags carriers is a coupling between a scalar
that is even in the thickness and a density that is also even; the Peltier heat does the same in the other
direction. There is no way to assign a single parity to the pair, so the two families mix. Everything else
in the model respects the pairing: the piezoelectric constants, the pyroelectric constant, the thermal
stress coefficients and the size operators all keep it. This is a stricter statement than the one we
started with, which asked for a polar coupling as well; the algebra shows that the polar couplings already
respect the pairing and only the two thermoelectric coefficients break it.

\subsection{What the parity means for the boundary rows}

When \eqref{eq:parity} holds, the five amplitudes split into two sets, one odd and one even under the sign
vector, and each set satisfies a smaller problem of its own. The two sets then give two smaller boundary
matrices, and the full determinant factorises into the product of the two. When it fails, there is one
$5\times5$ problem and one boundary matrix, and the symmetric and antisymmetric families are not separately
defined. Both cases are assembled in Section~\ref{sec:G}; you should be able to say which case your chosen
material falls in before you start any algebra.

\section{Calculation E: amplitude ratios and partial waves}\label{sec:E}

For each root $\eta_m$ of \eqref{eq:char} the matrix $\mathbf M(\eta_m)$ is singular, so the amplitudes are
not all independent: they are proportional to any nonzero column of the adjugate. We use the cofactor
formula, which needs no division and therefore survives a repeated or a nearly repeated root.

\subsection{Cofactors}

\begin{worked}
Let $\mathbf A(\eta)$ be the adjugate of $\mathbf M(\eta)$, so that $\mathbf M\mathbf A=\det(\mathbf M)\mathbf I$.
The entry $A_{ij}$ is $(-1)^{i+j}$ times the determinant of the $4\times4$ matrix obtained by deleting row
$j$ and column $i$ from $\mathbf M$. At a root the product $\mathbf M\mathbf A$ vanishes, so every column of
$\mathbf A$ is a null vector. Take the column with the largest entries and divide by its largest entry; the
ratios are then finite and one of them equals one by construction.
\end{worked}

\noindent For the root $\eta_m$ write the amplitudes as
\begin{equation}\label{eq:ratio}
\bigl(\bx_m,\bz_m,\bp_m,\bn_m,\bt_m\bigr)=\bt_m\bigl(\rho^{(1)}_m,\rho^{(2)}_m,\rho^{(3)}_m,\rho^{(4)}_m,1\bigr),
\end{equation}
where the three ratios are read from the fifth row of the adjugate, and where $\bt_m$ is the arbitrary
amplitude of the temperature part of that partial wave. The choice of the fifth entry as the normalising one
is a convention; it fails only for a root whose temperature amplitude vanishes, and in that case you
normalise by the largest of the five entries instead.

\subsection{The general solution through the thickness}

\begin{worked}
Each root contributes $\exp(\eta_m z)$ times its amplitude vector. Summing over the $N$ roots of
\eqref{eq:char}, and multiplying by the propagation factor,
\[
u_x=\sum_{m=1}^{N}a_m\rho^{(1)}_m e^{\eta_m z},\quad
u_z=\sum_{m=1}^{N}a_m\rho^{(2)}_m e^{\eta_m z},\quad
\varphi=\sum_{m=1}^{N}a_m\rho^{(3)}_m e^{\eta_m z},
\]
\[
\Delta n=\sum_{m=1}^{N}a_m\rho^{(4)}_m e^{\eta_m z},\quad
\theta=\sum_{m=1}^{N}a_m e^{\eta_m z},
\]
all multiplied by $\exp[\ii(kx-\omega t)]$, with $N=\deg_\eta\det\mathbf M$ and with the free constants
$a_m$ carrying the dimension of temperature.
\end{worked}

\noindent So the general solution of the field equations that is bounded in the plate is a sum of $N$
partial waves,
\begin{equation}\label{eq:general}
\bigl(u_x,u_z,\varphi,\Delta n,\theta\bigr)
=\sum_{m=1}^{N}a_m\bigl(\rho^{(1)}_m,\rho^{(2)}_m,\rho^{(3)}_m,\rho^{(4)}_m,1\bigr)
e^{\eta_m z}e^{\ii(kx-\omega t)},
\end{equation}
with $N=14$ when $l_g\ne0$ and $N=10$ when $l_g=0$. The number $N$ is exactly the number of boundary
conditions we are allowed to impose, seven per face when $l_g\ne0$ and five per face when $l_g=0$; that is
the content of Section~\ref{sec:F}, and it is why the boundary problem is square. If the number of roots and
the number of conditions disagreed, the problem would be over- or under-determined and no dispersion
relation would follow.

\subsection{One minor written out}

The cofactor formula is abstract until you do one, so here is one. The entry $A_{51}$ of the adjugate is the
cofactor of $M_{15}$, that is $(-1)^{5+1}=+1$ times the determinant of the matrix left after deleting row 1
and column 5:
\begin{equation}\label{eq:minor}
A_{51}=\begin{vmatrix}
M_{21}&M_{22}&M_{23}&M_{24}\\
M_{31}&M_{32}&M_{33}&M_{34}\\
M_{41}&M_{42}&M_{43}&M_{44}\\
M_{51}&M_{52}&M_{53}&M_{54}
\end{vmatrix}.
\end{equation}

\begin{worked}
Put the entries of Section~\ref{sec:C} into \eqref{eq:minor}. Row 1 of the minor is the second row of
$\mathbf M$ with its fifth entry deleted: $-\ii k\eta\Gc(c_{13}+c_{44})$, $-\Nc\rho\omega^2+\Gc(c_{44}k^2-c_{33}\eta^2)$,
$e_{15}k^2-e_{33}\eta^2$, $0$. Row 2 is the Gauss row: $\ii k\eta(e_{15}+e_{31})$, $e_{33}\eta^2-e_{15}k^2$,
$\epsilon_{11}k^2-\epsilon_{33}\eta^2$, $q$. Row 3 is the carrier row: $0$, $0$, $b_nn_0L$,
$-\ii\omega-D_nL-\ii kv_d+1/\tau_R$. Row 4 is the heat row: $fT_0\omega k\beta_1$,
$-\ii\omega fT_0\beta_3\eta$, $\ii\omega fT_0p_3\eta$, $f(-\ii\omega q\Pth+E_g/\tau_R)$. Expanding along the
first column of the minor leaves two terms, because the carrier row has a zero there. Each of those two
terms is a $3\times3$ determinant whose entries you already have, so the whole cofactor is arithmetic on
known quantities once $\eta_m$ and $k$ are fixed.
\end{worked}

\noindent The other four entries of row 5 of the adjugate follow the same pattern with a different column
deleted, and the amplitude ratios of \eqref{eq:ratio} are their ratios after normalising by the largest. If
you check one thing in this section, check that $A_{51}$ reduces to the elastic expression when the electric,
carrier and thermal couplings are switched off: then the minor becomes triangular and its determinant is the
product of the two mechanical entries, which you can verify by hand.

\subsection{Why the count is not a coincidence}

\begin{worked}
Look at the order of the field equations. Four of the five equations are of second order in $z$ when
$l_g=0$: momentum twice, Gauss once, carriers once, heat once, but Gauss carries the second derivative of
the potential only, so five second-order equations would suggest ten conditions per face pair, that is
five per face. The fifth-order character comes from the size operators: each factor
$(1-l_g^2\partial_z^2)$ in the elastic entries adds two orders, and it appears in the determinant in the
product $\Gc^2$, so four extra roots appear and four extra conditions, two per face.
\end{worked}

\noindent One rule matches all of this: each factor of $(1-l^2\partial_z^2)$ in the field equations adds one
root pair and needs one extra condition per face. The gradient length enters squared in the determinant, so
it adds two pairs. The nonlocal length enters only in the inertia and adds no root pair, which is why it
also adds no boundary condition. This is a useful check on any matrix you build: count the factors of the
gradient operator in the entries, double it, and add it to the local count of ten.

\section{Calculation F: the boundary conditions}\label{sec:F}

The faces $z=\pm h$ are free, electrically open, free of carrier flux and thermally insulating. In a
gradient theory the first two of these need care, because the variation of the energy at a free surface
carries a surface term with a gradient of the displacement. We follow the treatment of Charalambopoulos,
Douka and Mavratzas \cite[Sec.~2, Eq.~(14)]{Charalambopoulos2019}, who write the traction condition for a
gradient elastic solid with the surface gradient of the double traction. The form we use is the following.

\subsection{The seven conditions}

At each face, with outward normal $\mathbf n=(0,0,\pm1)$ and with $z_0=\pm h$:

\begin{enumerate}
\item[(b1)] the normal traction vanishes, with the surface-gradient term:
$\bigl[\sigma_{zz}-\partial_x R_{xz}\bigr]_{z=z_0}=0$, where $R_{ij}=n_k\mu_{kijl}n_l$ is the double
traction of the gradient theory;
\item[(b2)] the shear traction vanishes, with its surface-gradient term:
$\bigl[\sigma_{xz}-\partial_x R_{xx}\bigr]_{z=z_0}=0$;
\item[(b3)] the double traction normal to the face vanishes: $\bigl[R_{zz}\bigr]_{z=z_0}=0$;
\item[(b4)] the double traction in the plane of the face vanishes: $\bigl[R_{zx}\bigr]_{z=z_0}=0$;
\item[(b5)] the electric condition is open circuit: $\bigl[D_z\bigr]_{z=z_0}=0$;
\item[(b6)] no carriers leave or enter: $\bigl[\Gamma_z\bigr]_{z=z_0}=0$, where $\Gamma_z$ is the normal
component of the carrier flux;
\item[(b7)] the face is thermally insulating: $\bigl[\partial_z\theta\bigr]_{z=z_0}=0$.
\end{enumerate}

\noindent Conditions (b1) and (b2) are the standard free-surface conditions of the gradient theory with the
divergence of the surface double traction added; (b3) and (b4) are the accompanying double-traction
conditions, which in the language of \cite[condition b(i)]{Charalambopoulos2019} state that the normal
double force vanishes. Conditions (b5), (b6) and (b7) are our modelling choices and carry no source: open
circuit, no carrier flux, adiabatic face. Each is a choice, and each has an alternative; we list the
alternatives in the self-test section.

The word ``ours'' must be understood precisely here. Equations (b1) and (b2) with the surface gradient are
taken from \cite[Sec.~2, Eq.~(14)]{Charalambopoulos2019} for the purely elastic case; extending the same
structure to the piezoelectric and thermal parts, that is letting $\sigma$ and $R$ carry the electric and
thermal contributions, is our step. If you want a conservative version, keep the surface-gradient terms in
the elastic part only and drop them from the electric and thermal parts; the matrix changes in four entries
and the dispersion curves move by a small amount.

\subsection{The nonlocal operator, inverted root by root}

The constitutive law \eqref{eq:const} is written with the gradient operator applied to the strain. In the
integral form of NSGT the same physics is written with a kernel, and the differential form we use is
equivalent to it only when the constitutive boundary conditions of \cite[Eq.~(19)]{Barretta2018} are
satisfied. We do not impose those conditions as extra rows; instead we follow the route the supervisor
fixed: the nonlocal operator is inverted root by root, so that each partial wave carries the factor
$1/\Nc(\eta_m)$ in its stress.

\begin{worked}
The reason is one line of algebra. Suppose the stress amplitude of a partial wave is $\hat\sigma_m$, and the
constitutive law at that root reads $\Nc(\eta_m)\hat\sigma_m=\hat\sigma^{\rm loc}_m$, where the local part is
the right side of \eqref{eq:sxx} without the operator. Then $\hat\sigma_m=\hat\sigma^{\rm loc}_m/\Nc(\eta_m)$.
In the boundary rows we need the stress itself, so we divide each partial wave's stress by its own
$\Nc(\eta_m)$ rather than by a common operator. Dividing by a common $\Nc(\eta)$ would be wrong: the
operator does not commute with the sum over roots when the roots are complex.
\end{worked}

\noindent Result: in every boundary row, the stress contribution of partial wave $m$ is divided by
$\Nc(\eta_m)$, and the factor is inside the column, so it multiplies the whole column $m$.

\subsection{From conditions to row functions}

\begin{worked}
Take (b1) at the face $z=+h$ and put the partial wave $m$ in. The stress is $\hat\sigma_{zz}$ of
\eqref{eq:szz} with the amplitudes of that wave, divided by $\Nc(\eta_m)$; the surface-gradient term is
$-\partial_xR_{xz}$, and $\partial_x\to\ii k$ acting on the exponential, so the term is
$-\ii kR_{xz}$. The double traction for a normal along $z$ is $R_{xz}=\mu_{xzkl}n_kn_l$; in the gradient
theory with one internal length $l_g$ the only non-zero components for our class are proportional to
$l_g^2c_{44}$ acting on the gradient of $\eps_{xz}$, so within the differential law the term can be written
as $-l_g^2c_{44}\,\partial_x\nabla^2\eps_{xz}$, and under the plane wave this is
$-\ii k\,l_g^2c_{44}(\eta_m^2-k^2)\,\hat\eps_{xz}$. Every factor here is known once $\eta_m$ is known, so
the row is a linear function of the amplitudes.
\end{worked}

\noindent This gives the following seven row functions for the face $z=+h$, each to be evaluated with the
amplitudes of partial wave $m$ and then multiplied by $\exp(\eta_mh)$:

\begin{align}
b_1&=\frac{1}{\Nc(\eta_m)}\Bigl[\hat\sigma_{zz}-\ii k\,l_g^2c_{44}(\eta_m^2-k^2)\hat\eps_{xz}\Bigr],
\label{eq:b1}\\
b_2&=\frac{1}{\Nc(\eta_m)}\Bigl[\hat\sigma_{xz}-\ii k\,l_g^2c_{11}(\eta_m^2-k^2)\hat\eps_{xx}\Bigr],
\label{eq:b2}\\
b_3&=\frac{1}{\Nc(\eta_m)}\,l_g^2c_{33}(\eta_m^2-k^2)\hat\eps_{zz},
\label{eq:b3}\\
b_4&=\frac{1}{\Nc(\eta_m)}\,l_g^2c_{44}(\eta_m^2-k^2)\hat\eps_{xz},
\label{eq:b4}\\
b_5&=\hat D_z ,\qquad
b_6=\hat\Gamma_z ,\qquad
b_7=\eta_m\hat\theta .
\label{eq:b567}
\end{align}

The seven conditions (b1)--(b7) are the rows; the boundary matrix has seven rows per face.

\subsection{The carrier flux}

\begin{worked}
The electron current density in the drift-diffusion picture is
$J_z=q\bigl(b_nn_0E_z+D_n\partial_z\Delta n+b_nn_0S_n\partial_z\theta\bigr)$, which is the current whose
divergence appears in \eqref{eq:carr} once the thermodiffusion term is written with $S_n$ instead of $\bth$.
With $E_z=-\eta\bp$ and $\partial_z\to\eta$,
\[
\hat J_z=q\Bigl(-b_nn_0\eta\bp+D_n\eta\bn+b_nn_0S_n\eta\bt\Bigr).
\]
The condition of zero flux is $\hat J_z=0$; dividing by the common factor $q\eta$ leaves the row
$b_6=-b_nn_0\bp+D_n\bn+b_nn_0S_n\bt$. The factor $\eta$ is not zero at a root in general, so the division is
allowed; if $\eta_m=0$ the row is identically zero and must be replaced by the second-order condition,
which happens only for a degenerate root.
\end{worked}

\noindent The carrier row is
\begin{equation}\label{eq:b6}
b_6=-b_nn_0\bp+D_n\bn+b_nn_0S_n\bt .
\end{equation}

\subsection{The two faces}

\begin{worked}
At $z=-h$ the same seven functions appear, but with two changes. The amplitudes of a partial wave carry the
factor $\exp(-\eta_mh)$ instead of $\exp(\eta_mh)$. The outward normal is $(0,0,-1)$, so the sign of the
normal components flips: the traction $\hat\sigma_{zz}$ enters with a minus, the shear $\hat\sigma_{xz}$ with
a minus, the induction $D_z$ with a minus and the flux with a minus, while the double tractions, which are
quadratic in the normal, do not change sign, and the temperature gradient $\eta_m\hat\theta$ does change
sign because $\partial_z$ is odd.
\end{worked}

\noindent Writing $c_m=\exp(\eta_mh)$ and $d_m=\exp(-\eta_mh)$, the fourteen rows are
\begin{equation}\label{eq:rows}
\begin{aligned}
&b_1c_m,\ b_2c_m,\ b_3c_m,\ b_4c_m,\ b_5c_m,\ b_6c_m,\ b_7c_m &&\text{at } z=+h,\\
-&b_1d_m,\ -b_2d_m,\ b_3d_m,\ b_4d_m,\ -b_5d_m,\ -b_6d_m,\ -b_7d_m &&\text{at } z=-h .
\end{aligned}
\end{equation}

\section{Calculation G: the boundary matrix and the dispersion relation}\label{sec:G}

\subsection{Assembly}

Number the roots $\eta_1,\dots,\eta_N$ in any fixed order, with $N$ as in \eqref{eq:deg}. The boundary
matrix $\mathbf B$ has one row per condition of \eqref{eq:rows} and one column per root. Row $i$, column $m$
is the row function $b_i$ of \eqref{eq:b1}--\eqref{eq:b6} evaluated with the amplitudes of root $\eta_m$,
multiplied by the factor in \eqref{eq:rows}.

\begin{worked}
The amplitude ratios in $b_i$ are the cofactors of Section~\ref{sec:E}, so each entry of $\mathbf B$ is a
product of three known factors: the material constant of that row, the amplitude ratio of that column, and
the exponential factor of that face. Nothing in an entry depends on $k$ except through the roots $\eta_m$
and through the explicit $k$ in the row functions; the whole $k$-dependence sits inside $\mathbf B$.
\end{worked}

\noindent The dispersion relation is
\begin{equation}\label{eq:disp}
\det\mathbf B(k,\omega)=0 ,
\end{equation}
a $14\times14$ determinant when $l_g\ne0$ and a $10\times10$ determinant when $l_g=0$. Note that when
$l_g=0$ the rows (b3) and (b4) are identically zero, since every entry carries a factor $l_g^2$; they must be
deleted from the matrix, together with the four roots that the vanishing gradient operator removes. That is
how the size of the matrix falls from fourteen to ten, and it is the reason the two cases must be written
separately in a numerical treatment.

\subsection{What is solved for}

\begin{worked}
The frequency $\omega$ is real and given, because an experiment drives the plate at a known frequency. The
wave number is complex, $k=k_r+\ii k_i$. For each real $\omega$ the equation \eqref{eq:disp} is one complex
equation in two real unknowns, so its solutions are isolated pairs. The real part gives the phase velocity
$\omega/k_r$ and the imaginary part gives the attenuation; a positive $k_i$ means the wave decays along
$+x$, which is the sign convention of this file.
\end{worked}

\noindent Solving \eqref{eq:disp} for real $\omega$ means finding the zeros of a determinant whose entries
are themselves roots of \eqref{eq:char}. There is no closed form; the practical route is a search in the
complex $k$ plane started from the uncoupled elastic solution, and the search must be watched because two
branches can approach each other closely when the carrier relaxation crosses the wave frequency.

\subsection{The column scaling}

One technical remark belongs here, because it is the only place where the arithmetic of a computer can
mislead you. In \eqref{eq:rows} the factors $\exp(\eta_mh)$ can be enormous when $\eta_m$ has a large real
part, which happens for the roots created by the size operator. Their products with the amplitude ratios
then differ by hundreds of orders of magnitude between columns, and a determinant evaluated with such
columns loses all its digits. The cure is to multiply column $m$ by $\exp(-|\mathrm{Re}\,\eta_m|h)$ before
any evaluation. A common factor per column does not change the zeros of the determinant, so the dispersion
relation is untouched, but the arithmetic becomes well conditioned.

\subsection{What happens when the parity holds}

When \eqref{eq:parity} holds, the determinant of Section~\ref{sec:D} is even and the roots come in pairs
$\pm\eta$. The five amplitudes split into an odd set and an even set. Assembling the seven conditions in the
same two groups then gives two matrices, one for each family, and
\begin{equation}\label{eq:factor}
\det\mathbf B=\det\mathbf B_{\rm S}\cdot\det\mathbf B_{\rm A},
\end{equation}
where the labels S and A stand for the symmetric and antisymmetric families of the classical plate. Each of
the two smaller matrices has seven rows and seven columns. This factorisation is not assumed anywhere; it is
a consequence of the parity and it must appear by itself when you build the matrices. If it does not appear,
one of your rows has the wrong parity, and the likeliest place is a sign in \eqref{eq:b1} or in the flux
\eqref{eq:b6}.

\section{Calculation H: hand limits}\label{sec:H}

\subsection{All couplings off: the Rayleigh--Lamb relation}

\begin{worked}
Set $e_{31}=e_{33}=e_{15}=0$, $p_3=0$, $\beta_1=\beta_3=0$, $\bth=\Pth=0$, $\tau_R\to\infty$, and let the
carrier and thermal fields drop out. Rows 3, 4 and 5 then decouple from rows 1 and 2, and $\det\mathbf M$ is
the product of the $2\times2$ mechanical determinant and the three scalar factors. Keep only the mechanical
part and set $l_g=l_{\rm nl}=0$. The mechanical determinant reads
\[
\begin{vmatrix}
-\rho\omega^2+c_{11}k^2-c_{44}\eta^2 & -\ii k\eta(c_{13}+c_{44})\\
-\ii k\eta(c_{13}+c_{44}) & -\rho\omega^2+c_{44}k^2-c_{33}\eta^2
\end{vmatrix}=0 ,
\]
which is the classical P--SV condition for a transversely isotropic plate. Its four roots are
$\eta=\pm\ii p$ and $\eta=\pm\ii q$ with $p^2=\omega^2/c_L^2-k^2$ and $q^2=\omega^2/c_T^2-k^2$, where
$c_L$ and $c_T$ are the longitudinal and transverse speeds built from $c_{33},c_{13},c_{44}$ along the
axis. For an isotropic solid, $c_{11}=c_{33}=\lambda+2\mu$, $c_{13}=\lambda$, $c_{44}=\mu$ and
$c_L^2=(\lambda+2\mu)/\rho$, $c_T^2=\mu/\rho$.
\end{worked}

\noindent The four boundary rows that survive are the two tractions at each face, and the determinant of
that $4\times4$ system reduces by the standard trigonometric elimination to the two Rayleigh--Lamb
equations,
\begin{equation}\label{eq:RL}
\frac{\tan(qh)}{\tan(ph)}=-\frac{4k^2pq}{(k^2-q^2)^2}\quad\text{(symmetric)},\qquad
\frac{\tan(qh)}{\tan(ph)}=-\frac{(k^2-q^2)^2}{4k^2pq}\quad\text{(antisymmetric)} .
\end{equation}

\begin{worked}
The elimination is worth doing once by hand. The $4\times4$ determinant has the structure of a difference of
two $2\times2$ blocks because the symmetric combination pairs $\cos(ph)$ with $\sin(qh)$ and the
antisymmetric pairs $\sin(ph)$ with $\cos(qh)$. Add and subtract the rows of the two faces; the sum of the
two traction rows at $z=\pm h$ keeps the even functions and the difference keeps the odd ones. What remains
is a $2\times2$ determinant in each case, and setting it to zero gives \eqref{eq:RL} after dividing by
$\sin(ph)\sin(qh)$ and using $\cos/\sin=\cot$.
\end{worked}

\noindent The two families in \eqref{eq:RL} are the S and A families. Their long-wavelength limits are the
classical ones: the lowest symmetric branch tends to the plate speed $\sqrt{E/(\rho(1-\nu^2))}$ as the
frequency--thickness product goes to zero, the lowest antisymmetric branch tends to zero as a bending wave,
and both tend to the Rayleigh speed of the half space as the product becomes large. A numerical treatment
that does not reproduce these three statements has a mistake somewhere, and the mistake is usually a sign
in one of the stress components \eqref{eq:sxx}--\eqref{eq:sxz}.

\subsection{Thin plate, open circuit: the stiffened speed}

\begin{worked}
Take the frequency--thickness product small, so that the wavelength is long compared with $2h$ and the
fields are almost uniform across the thickness. Keep the piezoelectric coupling and set the size lengths to
zero. The faces are free and open, so $\sigma_{zz}=\sigma_{xz}=D_z=0$ at both faces, and the plate carries
only $\eps_{xx}$ and $E_z$. Eliminating $E_z$ from $D_z=0$ gives
$E_z=(e_{31}\eps_{xx}+p_3\theta)/\epsilon_{33}$, and putting this into $\sigma_{xx}$ gives an effective
stiffness
\[
\bar c_{11}=c_{11}+\frac{e_{31}^2}{\epsilon_{33}} ,
\]
which is larger than $c_{11}$: the open circuit makes the plate stiffer than the short circuit, where $E_z$
is zero and the stiffness is $c_{11}$ itself. With the density unchanged, the extensional speed of the thin
plate becomes
\[
c_{\rm open}=\sqrt{\frac{\bar c_{11}}{\rho}}\quad(\text{stiffened}),\qquad
c_{\rm short}=\sqrt{\frac{c_{11}}{\rho}}\quad(\text{short circuit}).
\]
If the faces are not traction free in the $z$ direction, the elimination also brings in $c_{13}$ and $c_{33}$
and the combination is $\bar c_{11}=c_{11}-c_{13}^2/c_{33}+(e_{31}-c_{13}e_{33}/c_{33})^2/(\epsilon_{33}+e_{33}^2/c_{33})$;
for the plate we work with, the first form is the one to compare against.
\end{worked}

\noindent The limit is a statement about stiffness and not about numbers: the open-circuit plate is stiffer
than the short-circuit plate, the difference grows with the square of the piezoelectric constant and falls
with the permittivity, and it is the same difference that appears as the electromechanical coupling factor in
the literature. A check that the coupled matrix reproduces this ordering at small frequency--thickness is the
first test of the piezoelectric block of \eqref{eq:M}, because it tests the Gauss row and the open-circuit
condition together.

\subsection{Bulk limit: thermoelastic damping}

\begin{worked}
Let $k$ be large compared with $\omega/c_L$, so that the wave is a plane wave in an unbounded medium as far
as the heat is concerned, and keep only the thermal coupling to the dilatation. The heat equation
\eqref{eq:g5} with $f=1$ reads
$-\ii\omega\rho c_p\bt+T_0\omega k\beta_1\bx-\kB L\bt=0$, and the mechanical equation carries $-\beta_1\bt$
as a stress. Eliminating $\bt$ between the two gives a complex correction to the elastic modulus whose
imaginary part, divided by the real part, is
\[
\frac{k_i}{k_r}=\frac{\Delta_E}{2}\,\frac{\omega\tau}{1+(\omega\tau)^2},
\]
with $\Delta_E$ the dimensionless thermal expansion strength and $\tau$ the thermal relaxation time.
\end{worked}

\noindent This is the classical thermoelastic damping formula of the Fourier theory. Two things happen when
the MGT operator is kept. The denominator becomes
$1+\omega^2(\tau_T-\tau_q)^2$ in place of $1+(\omega\tau)^2$, so the peak moves and its height changes; and
the numerator gains the factor $f$ from the relaxation. The peak of the damping is the classical
thermoelastic peak, at $\omega\tau=1$ in the Fourier case.

\subsection{Bulk limit: the Hutson--White condition}

\begin{worked}
The carrier equation \eqref{eq:g4} at long wavelength, with the field term and the drift term dropped,
reads $(-\ii\omega-D_nL+1/\tau_R)\bn=0$, so a disturbance decays at the rate $D_n|L|$ and is swept by the
wave at the rate $\omega$. The two rates are equal when $D_nk^2=\omega$, that is when
$k^2=\omega/D_n$. Using the Einstein relation $D_n=b_nk_BT_0/q$ and the definition of the Debye length of
the carriers, the equality between the dielectric relaxation rate and the wave frequency reads
\[
n_0=\frac{\epsilon_{11}\omega}{q\,b_n},
\]
which is the condition for the acoustoelectric interaction to be strongest: the carriers have just enough
time to follow the wave.
\end{worked}

\noindent This condition is the one to test against a scan in the equilibrium density: at fixed frequency,
the attenuation as a function of $n_0$ has a maximum at the value above. It is a check on row 4 of
\eqref{eq:M} and on the carrier row \eqref{eq:b6} of the boundary matrix, and it is the reason the carrier
density must be kept as a variable in any numerical study rather than set to a fixed value.

\subsection{Other limits worth doing by hand}

\begin{worked}
Three more reductions cost one line each. Setting $l_g=l_{\rm nl}=0$ returns the classical plate and
\eqref{eq:RL} exactly. Setting $\tau_q=0$ returns the Fourier heat equation and the classical thermoelastic
damping. Setting $\bth=\Pth=0$ restores the parity of \eqref{eq:parity} and splits the matrix into the two
families; if the split does not appear, the error is in row 4 or row 5. Each of these is a switch in one
entry, so each is a check you can run in an afternoon.
\end{worked}

\section{Self-test list}\label{sec:tests}

These are the checks to run before you show the draft to anyone. Each is stated so that you can answer yes
or no.

\begin{enumerate}
\item Units. Every entry of \eqref{eq:row1}--\eqref{eq:row5} must have the same unit within its row, because
the equation is a sum. Check the units of $M_{15}=\ii k\beta_1$ (Pa/K per m) against $M_{11}$ (Pa/m) and
against $M_{13}$ (C/m$^3$ times V/m). The table in Section~\ref{sec:not} gives the units.
\item Symmetry under $\eta\to-\eta$. With $\bth=\Pth=0$ the determinant must be even, and with either of them
on it must not. Test both.
\item Degree. With $l_g\ne0$ the root count must be fourteen, with $l_g=0$ it must be ten. If you get a
different number, look at the highest powers in rows 1 and 2.
\item The classical limit. With all couplings off, both families of \eqref{eq:RL} must come out, with the
three long- and short-wavelength limits named in Section~\ref{sec:H}.
\item The open-circuit ordering. At small frequency--thickness the open-circuit speed must exceed the
short-circuit speed, with the ratio set by $e_{31}^2/(\epsilon_{33}c_{11})$.
\item The thermal peak. The attenuation against frequency must have a maximum near $\omega\tau_T=1$ in the
Fourier limit, and the maximum must move when $\tau_q$ is switched on.
\item The carrier peak. The attenuation against $n_0$ must have a maximum at the condition of
Section~\ref{sec:H}.
\item The size effect. Increasing $l_g$ must stiffen the plate and increasing $l_{\rm nl}$ must soften it.
If both do the same thing, the signs of $\Nc$ and $\Gc$ are the same in your code, which is wrong.
\item The parity of the boundary rows. Each of the seven rows of \eqref{eq:b1}--\eqref{eq:b6} must have a
definite parity under $\eta\to-\eta$ when \eqref{eq:parity} holds. A row that mixes the two parities is a
row with a term in the wrong place.
\item The alternatives. Rewrite (b5) as closed circuit, (b6) as a fixed carrier density and (b7) as an
isothermal face, one at a time, and see that each changes the matrix in one row only.
\end{enumerate}

\section{Work plan}

\begin{enumerate}
\item Week 1. Read Lim, Zhang and Reddy, Section 2, equations (17) to (19), and Barretta and Marotti de
Sciarra, equations (18) and (19), side by side. Write the differential law and the constitutive boundary
conditions in your own words. Confirm the placement of the operators in \eqref{eq:M}.
\item Week 2. Read the block of equations in Zhu et al.\ and attach the exact number to each of our
equations. Read Quintanilla's equations (2.5) and (3.1)--(3.5) and confirm the operator in \eqref{eq:heat}.
\item Week 3. Redo Calculations A to D on paper without looking at this file. Compare, and mark every place
where you differ.
\item Week 4. Redo Calculations E to G. Build the matrix for a chosen material, and check the ten self-tests
in order.
\item Week 5. Sensitivity study: vary each of the two length parameters and the two time constants one at a
time, and write down which of the ten tests moves and by how much. This is the part that becomes the results
section of the paper.
\item Week 6. Write the paper: the introduction from Section~2, the model from Sections~3 to 7, the
reduction from Sections~8 to 12, and the limits from Section~13. Put the worked blocks back only where your
supervisor asks for them.
\end{enumerate}

\bibliographystyle{unsrt}
\bibliography{Lamb_Sumit}

\end{document}
